\section{调试与可视化}

\subsection{Prompt 级别的仪表盘}

在部署后，wrap-around instrumentation 包括：prompt signature fingerprints、action entropy、KL drift trace、human override counter。将这些信号写入 time-series dashboard，可以快速定位哪类 prompt 开始触发 drift。

\begin{table}[H]
\centering
\begin{tabular}{p{0.25\textwidth} p{0.65\textwidth}}
\toprule
指标 & 说明 \\
\midrule
Prompt signature & hash + metadata，用于重新播放与 reproduce \\
Entropy & action entropy drop 可能预示 policy collapse \\
KL drift & 当前 policy vs. reference distribution \\
Override count & human 改写提示的次数，衡量 automation trust \\
\bottomrule
\end{tabular}
\caption{Prompt 级别仪表盘的核心信号}
\end{table}

\begin{importantbox}{指标要与 runbook 绑定}
每个 signal 的 threshold 要写入 runbook，并指定 owner。如果超限，就要自动通知对应的人（ops、alignment、legal），而不是靠 slack 注释。
\end{importantbox}

\subsection{Rollout 可视化}

\begin{figure}[H]
\centering
\includegraphics[width=0.8\textwidth]{slide-02-04.jpg}
\caption{Slide 02-04 把 rollout trajectory、human override 与 drift indicator 同步显示，形成统一的可视化面板。}
\end{figure}

可视化面板允许 engineer 在一分钟之内判断：1) 该 rollout 是否偏离 expert manifold；2) human override 是否集中在某类 prompt；3) KL/entropy 曲线是否同步下降。把这些图表 embed 到 internal dashboard，用作 gate review。

\begin{knowledgebox}{可视化的诊断价值}
当 rollout 出现 drift，visualization 让团队不需翻日志即可看到具体 prompt、KL 路径和 override 位置，大幅缩短 incident response 时间。
\end{knowledgebox}

\subsection{本章小结}

调试与可视化把抽象的 drift 信号具体化：Prompt dashboard 提供定量监控，Rollout visualization 则提供触发点的 qualitative insight，两者结合，才能在下一次 release 前完成根因分析。

